\chapter{Implementation Details} \label{cap:impl}

\section{Meeting Workflow}

The core meeting workflow is implemented in \texttt{MeetingService}. When a user proposes a meeting, the service creates a \texttt{Meeting} entity and immediately adds the organizer as an accepted participant. Each invited username is trimmed, empty values are ignored, duplicate invitees are skipped, and unknown invitees cause the operation to fail.

Invitation responses are represented by the \texttt{InviteStatus} enum. The service only accepts \texttt{ACCEPTED} or \texttt{DECLINED} as valid responses. Calendar queries exclude declined invitations for the invitee, which means a declined meeting no longer blocks that user's calendar, while the organizer can still see the meeting they created.

\section{Discovery and Calendar Export}

Events copied from discovery providers are transformed into normal meetings with a single accepted participant. If the provider does not supply an end time, the application assigns a default duration of two hours. This keeps copied events simple and allows them to appear immediately in the user's calendar and iCalendar feed.

The iCalendar renderer writes a calendar document with one event entry per meeting. Meeting confirmation state is mapped to calendar status: confirmed meetings are exported as confirmed, while meetings with pending participants are exported as tentative. Basic escaping is applied to text fields before rendering.

\section{Validation and Error Handling}

\begin{itemize}
    \item \textbf{Meeting validation.} The service rejects missing titles, missing start or end times, and intervals where the end is not after the start. A null invitee list is treated as empty, while individual invitee names are trimmed and deduplicated.
    \item \textbf{User validation.} Registration trims usernames before checking uniqueness and saving the account. Blank usernames and blank passwords are rejected before password encoding or persistence.
    \item \textbf{Discovery validation.} Copied discovered events are converted into accepted single-user meetings. When a provider supplies an invalid end time, the service rejects the interval; when no end time is supplied, it applies the default two-hour duration.
    \item \textbf{Response validation.} The meeting response controller accepts only explicit \texttt{accept} and \texttt{decline} actions. Unknown response actions return a bad request instead of being interpreted as a decline.
\end{itemize}

This split keeps business validation close to the service methods while still letting the controllers provide appropriate HTTP behavior. The corresponding unit and integration tests exercise these cases directly.
